% 原文 p.112（後半）


\subsectionlink{まとめ}
本章の主要な論点と結果のいくつかをまとめよう。
\begin{itemize}
\item 量子状態は、古典物理学の状態と根本的に異なる。測定は一般に量子状態を根本的に変えるため、測定を追加して任意に拡張できる、物理量の値の一覧を単に記述するものではない。
\item コヒーレントな重ね合わせは量子力学の核心にあり、状態の古典的集団とは解釈できない。重ね合わせは個々の物理状態を記述する。実験室では、その存在と量子的性格を、典型的に干渉実験によって検証する。
\item 量子もつれは、系間の非局所的な量子的相関を記述し、デコヒーレンスの核心にある。それは、系間の確率関係という古典的概念を、はるかに超える。
\item 密度行列の概念は、デコヒーレンスの形式的記述で重要な役割を果たす。異なる三種類の密度行列を区別した。
\begin{itemize}
\item 純粋状態の密度行列は、純粋状態と直接同値な密度行列であり、系の完全に既知の物理状態を表す。
\item 混合状態の密度行列は、純粋状態の古典的確率分布、すなわち無知として解釈できる集団を記述する。
\item 縮約密度行列は、より大きな系の部分系に対して行いうる、すべての測定の統計を網羅的に記述する。デコヒーレンスの文脈では、環境ともつれた系の局所的測定統計を記述するために用いる。縮約密度行列は、関心のある系以外のすべての自由度について、非ユニタリなトレース演算を行って得る。それは非真正集団を表すので、混合状態の密度行列と解釈してはならない。
\end{itemize}
\item 測定問題、および、より一般的な量子から古典への移行の問題は、主要な三つの問題からなる。
\begin{itemize}
\item 優先基底の問題。何が経験上の優先される物理量を決めるのか。
% 原文 p.113
\hypertarget{srcpage:113}{}\label{srcpage:113}
\item 干渉が観測されない問題。なぜ干渉効果を観測するのは、それほど難しいのか。
\item 結果の問題。そもそも、なぜ測定には結果があるように見え、何が観測される特定の結果を選ぶのか。
\end{itemize}
本章で示し、本書の他の箇所、たとえば第8章でさらに論じるように、デコヒーレンスは最初の二問題を解決できる一方、第三の問題は、大半がデコヒーレンスの範囲外の解釈の事柄と、本質的に結び付いていると結論するのが妥当である。
\item Bohrの有名な「相補性原理」は、量子もつれの帰結として理解できる。関心のある系$\mathcal S$ともつれた第二の系$\mathcal S'$は、次の意味で、$\mathcal S$の情報を符号化できる。$\mathcal S$の重ね合わせ状態の成分と量子的に相関する、$\mathcal S'$の相対状態の重なりが小さいほど、$\mathcal S'$は成分をよりよく識別でき、$\mathcal S$だけを測定して成分間の干渉を観測するのは、難しくなる。
\item デコヒーレンスは、物理系が、第二の系$\mathcal S'$の役割を果たす環境に浸っているという考えに基づく。系と環境の量子もつれを通じて、環境は系の特定の状態を連続的に監視する。特定の状況で、どの状態が監視されるかは、具体的な系と環境の相互作用の形に依存する。
\item この環境による監視は、系の特定の状態についての情報が、環境に符号化されることを意味する。系が、それらの状態の重ね合わせで記述されれば、原理上だけにせよ、環境を測定して成分状態を識別できるため、系の局所的水準で、成分間の干渉の減衰が強制される。重ね合わせの存在を観測的に確認するには、系と環境の複合系を測定する必要があり、物理的に現実的な状況の大半では、実際上のすべての目的に対して不可能である。このように、制御されない環境との量子もつれによって、コヒーレンスは、実際上不可逆に、より大きな系と環境の組合せへ非局在化し、系にしかアクセスできない観測者には、実効的に利用不能となる。
\item このデコヒーレンス過程は、互いに関連する三つの帰結を持つ。
\begin{itemize}
\item 環境が監視する状態間の、コヒーレンスの局所的減衰。これは測定問題の一成分、すなわち干渉が観測されない問題を解決する。
\item 系の優先状態、すなわち指針状態の環境誘起超選択。これらは環境との量子もつれ、したがってデコヒーレンスに、最も影響されない状態である。
% 原文 p.114
\hypertarget{srcpage:114}{}\label{srcpage:114}
その頑健性により、実効的に古典的と見なせる。大域的な系と環境の力学を記述する全Hamiltonianが与えられたとき、これらの指針状態を決める方法を論じた。環境誘起超選択は、自然界で、重ね合わせではなく位置や電荷などの、優先される物理量が現れることに、非常に一般的な説明を与え、したがって優先基底の問題に、実効的な解決を与える。
\item 環境が超選択する状態についての情報が、環境に頑健かつ冗長に符号化され、観測者が環境の断片を捕らえることによって、系の情報を得られること。この方法は、通常、観測者が情報を集める仕方をよく反映し、観測者が系に与える擾乱を最小にする。
\end{itemize}
\item デコヒーレンスが、古典的類似物のない、純粋に量子力学的な効果であることを強調した。特に、デコヒーレンスは散逸を伴う必要がない。伴う場合には、通常、特に巨視的尺度で、散逸の時間尺度よりはるかに短い時間に起こる。また古典的雑音過程は、概念的にも観測的にも、デコヒーレンスと明確に区別しなければならない。
\end{itemize}
